\documentclass[12pt]{article}
\usepackage[lmargin=1cm, rmargin=2cm,tmargin=2cm,bmargin=2cm]{geometry}
\usepackage{amsmath,amssymb,amsthm,bm}
\usepackage{enumitem} % for structured lists
\usepackage{graphicx}

\begin{document}

\section*{Probability Quiz with Answers}

\subsection*{1. Probability Axioms and Basic Events}
\noindent
\textbf{1.1. Question}\\
You roll a fair six-sided die. Define events:
\[
  A = \{\text{the result is an odd number}\}, \quad
  B = \{\text{the result is } \ge 4\}.
\]
\begin{enumerate}[label=(\alph*)]
\item Which of the following sets describes $A$?  
  \textbf{(A)} $\{1,3,5\}$,\quad 
  (B) $\{2,4,6\}$,\quad 
  (C) $\{1,2,3\}$,\quad 
  (D) $\{4,5,6\}$.
\item What is $P(A)$? (Short numeric.)
\item Which of these is $P(A \cup B)$?  
  \textbf{(A)} $\tfrac{2}{6}$,\quad 
  (B) $\tfrac{3}{6}$,\quad
  (C) $\tfrac{4}{6}$,\quad 
  \textbf{(D)} $\tfrac{5}{6}.$
\item Are $A$ and $B$ mutually exclusive?  
  (A) Yes,\quad \textbf{(B)} No.
\end{enumerate}

\noindent
\textbf{1.1. Answer}
\begin{itemize}
\item (a) Correct set for $A$: \textbf{(A)} $\{1,3,5\}$.
\item (b) $P(A)=\tfrac{3}{6} = 0.50.$
\item (c) $A \cup B = \{1,3,4,5,6\}$, so $P(A\cup B)=\tfrac{5}{6}\approx 0.83.$ 
  Correct choice: \textbf{(D)}.
\item (d) $A$ and $B$ intersect at $\{5\}$, so they are not mutually exclusive. 
  Correct: \textbf{(B)} No.
\end{itemize}

\hrule

\subsection*{2. Counting Techniques}

\textbf{2.1. Question}\\
A group has $6$ men and $6$ women (12 total). A committee of $4$ people is chosen at random.
\begin{enumerate}[label=(\alph*)]
\item How many possible committees of $4$ are there in total? (Short numeric.)
\item What is the probability that the committee has exactly 2 men and 2 women?  
  \begin{enumerate}[label=\Alph*.]
  \item $\displaystyle \frac{\binom{6}{2}\,\binom{6}{2}}{\binom{12}{4}}$
  \item $\displaystyle \frac{\binom{6}{2}\,\binom{6}{2}}{\binom{12}{2}\,\binom{12}{2}}$
  \item $\displaystyle \frac{\binom{12}{2}}{\binom{6}{4}}$
  \item $\displaystyle \binom{6}{2}\,\binom{6}{2}\quad (\text{no division})$
  \end{enumerate}
\item What is the probability that at least 3 members are women?
  \begin{enumerate}[label=\Alph*.]
  \item $\displaystyle \frac{\binom{6}{3}\,\binom{6}{1} + \binom{6}{4}\,\binom{6}{0}}{\binom{12}{4}}$
  \item $\displaystyle \frac{\binom{6}{4}\,\binom{6}{4}}{\binom{12}{4}}$
  \item $\displaystyle \frac{\binom{12}{3}}{\binom{12}{4}}$
  \item $\displaystyle \binom{6}{3}\,\binom{6}{1} + \binom{6}{4}\,\binom{6}{0}\ (\text{no denominator})$
  \end{enumerate}
\end{enumerate}

\noindent
\textbf{2.1. Answer}
\begin{itemize}
\item (a) Total committees: $\binom{12}{4}=495.$
\item (b) Probability of exactly 2 men, 2 women:
  \[
    \frac{\binom{6}{2}\,\binom{6}{2}}{\binom{12}{4}} \quad\text{(Choice (A))}.
  \]
\item (c) Probability of at least 3 women:
  \[
    \frac{\binom{6}{3}\,\binom{6}{1} + \binom{6}{4}\,\binom{6}{0}}{\binom{12}{4}} 
    \quad\text{(Choice (A))}.
  \]
\end{itemize}

\textbf{2.2. Question}\\
In a simplified ``birthday'' model with $10$ equally likely birthdays, you have a group of $3$ people.  
What is the probability that all $3$ have different birthdays? (Short formula or fraction.)

\noindent
\textbf{2.2. Answer}\\
\[
  P(\text{all different}) \;=\; \frac{10 \times 9 \times 8}{10^3} 
  \;=\; 0.72.
\]

\hrule

\subsection*{3. Independence and Conditional Probability}

\textbf{3.1. Question}\\
$A$ and $B$ are independent with $P(A)=0.3, P(B)=0.2.$ Which is $P(A\cap B)$?
\begin{enumerate}[label=\Alph*.]
\item 0.06
\item 0.10
\item 0.50
\item Cannot be determined
\end{enumerate}

\noindent
\textbf{3.1. Answer}\\
For independent events, $P(A\cap B)=P(A)P(B)=0.3\times0.2=0.06.$  
Correct choice: \textbf{(A)}.

\medskip

\textbf{3.2. Question}\\
Two coins in a box:
\begin{itemize}
\item Coin 1 (fair): $P(H)=0.5$.
\item Coin 2 (possibly double-headed): $P(H)=1$.
\end{itemize}
Pick one coin at random and flip it once. It shows Heads ($H$). Let $D$ = ``double-headed coin.''

\begin{enumerate}[label=(\alph*)]
\item Which set of values is correct for $\{P(D), P(H\mid D), P(H\mid D^c)\}$?
  \begin{enumerate}[label=\Alph*.]
  \item $P(D)=\tfrac12,\;P(H\mid D)=1,\;P(H\mid D^c)=\tfrac12$
  \item $P(D)=\tfrac12,\;P(H\mid D)=\tfrac12,\;P(H\mid D^c)=1$
  \item $P(D)=1,\;P(H\mid D)=\tfrac12,\;P(H\mid D^c)=0$
  \item $P(D)=\tfrac12,\;P(H\mid D)=0,\;P(H\mid D^c)=\tfrac12$
  \end{enumerate}

\item What is $P(D \mid H)$? (Short numeric.)
\end{enumerate}

\noindent
\textbf{3.2. Answer}
\begin{itemize}
\item (a) Correct: \textbf{(A)} $P(D)=\tfrac12,\;P(H\mid D)=1,\;P(H\mid D^c)=0.5.$
\item (b) By Bayes' rule:
  \[
    P(D\mid H) 
    \;=\; \frac{1 \times \tfrac12}{\,1 \times \tfrac12 + \tfrac12 \times \tfrac12\,}
    \;=\;\frac{0.5}{0.5+0.25}
    \;=\;\frac{2}{3}\approx 0.6667.
  \]
\end{itemize}

\hrule

\subsection*{4. Bayes' Rule in Practice}

\textbf{4.1. Question}\\
A disease affects 1 in 500 people, so $P(\text{disease})=0.002$. A test has 90\% true positive rate and 5\% false positive rate. You test positive. Which is
\[
  P(\text{disease}\mid \text{test positive}) \;=\; ?
\]
\begin{enumerate}[label=\Alph*.]
\item $\displaystyle \frac{0.90 \times 0.002}{\,0.90 \times 0.002 + 0.05 \times 0.998\,}$
\item $\displaystyle \frac{0.05 \times 0.998}{\,0.90 \times 0.002\,}$
\item $0.90 + 0.002 - 0.05$
\item $\displaystyle \frac{0.90}{\,0.90 + 0.05\,}$
\end{enumerate}

\noindent
\textbf{4.1. Answer}\\
Bayes' formula indicates choice \textbf{(A)} is correct.

\medskip

\textbf{4.2. Question}\\
$P(W)=0.3$ for Candidate~A to win. Poll says ``high chance'' (HC) with probability $0.8$ if A wins, $0.15$ if A loses. Which formula is $P(W\mid HC)$?
\begin{enumerate}[label=\Alph*.]
\item $\displaystyle \frac{0.8\times 0.3}{\,0.8\times 0.3 + 0.15\times 0.7\,}$
\item $0.3 + 0.8 - 0.15$
\item $\displaystyle \frac{0.3 \times 0.15}{\,0.8\,}$
\item $0.8\times 0.3$
\end{enumerate}

\noindent
\textbf{4.2. Answer}\\
Correct: \textbf{(A)}.

\hrule

\subsection*{5. Random Variables and Distributions}

\textbf{5.1. Question}\\
Let $X\sim \mathrm{Binomial}(n=5,\,p=0.3).$
\begin{enumerate}[label=(\alph*)]
\item Which is the correct PMF?
  \begin{enumerate}[label=\Alph*.]
  \item $\binom{5}{k}\,(0.3)^k (0.7)^{5-k}$
  \item $\binom{5}{k}\,(0.5)^5$
  \item $\frac{5}{k}(0.3)^k$
  \item $k\,(0.3)$
  \end{enumerate}
\item $E[X]$? (Short numeric.)
\item $\mathrm{Var}(X)$? (Short numeric.)
\item Expression for $P(X\ge 3)$?
  \begin{enumerate}[label=\Alph*.]
  \item $\displaystyle \sum_{k=3}^{5}\binom{5}{k}(0.3)^k(0.7)^{5-k}$
  \item $\binom{5}{3}(0.3)^3(0.7)^2$
  \item $(0.3)^3$
  \item $1-\sum_{k=3}^{5}\binom{5}{k}(0.3)^k(0.7)^{5-k}$
  \end{enumerate}
\end{enumerate}

\noindent
\textbf{5.1. Answer}
\begin{itemize}
\item (a) Correct: \textbf{(A)}. 
\item (b) $E[X]=n\,p=5\times 0.3=1.5.$
\item (c) $\mathrm{Var}(X)=n\,p\,(1-p)=5\times 0.3\times0.7=1.05.$
\item (d) $P(X\ge 3)=\sum_{k=3}^5\binom{5}{k}(0.3)^k(0.7)^{5-k}.$ Choice \textbf{(A)}.
\end{itemize}

\medskip

\textbf{5.2. Question}\\
$Y\sim \mathrm{Uniform}(2,6).$
\begin{enumerate}[label=(\alph*)]
\item Which PDF is correct?
  \begin{enumerate}[label=\Alph*.]
  \item $f_Y(y)=\tfrac{1}{4}$ for $2\le y\le 6$
  \item $f_Y(y)=\tfrac{1}{6}$ for $2\le y\le 6$
  \item $f_Y(y)=4$ for $0\le y\le 1$
  \item $f_Y(y)=\tfrac{1}{6-2}$ for $y\le2$
  \end{enumerate}
\item What is $P(3.5 \le Y \le 5)$? (Short numeric.)
\item What is $E[Y]$? (Short numeric.)
\end{enumerate}

\noindent
\textbf{5.2. Answer}
\begin{itemize}
\item (a) Correct: \textbf{(A)} $f_Y(y)=\frac{1}{4}$ on $[2,6].$
\item (b) $P(3.5\le Y\le 5)=\frac{5-3.5}{6-2}=\frac{1.5}{4}=0.375.$
\item (c) $E[Y]=\frac{2+6}{2}=4.$
\end{itemize}

\hrule

\subsection*{6. Mean and Variance}

\textbf{6.1. Question}\\
$Z$ takes values $\{0,1,2\}$ with probabilities $0.5,0.3,0.2.$
\begin{enumerate}[label=(\alph*)]
\item $E[Z]$ is
  \begin{enumerate}[label=\Alph*.]
  \item $0.3$
  \item $0.7$
  \item $0.5$
  \item $1.0$
  \end{enumerate}
\item $E[Z^2]$ is
  \begin{enumerate}[label=\Alph*.]
  \item $0.5^2 + 0.3^2 + 0.2^2$
  \item $0.5\cdot0 + 0.3\cdot1 + 0.2\cdot2$
  \item $0^2\cdot0.5 + 1^2\cdot0.3 + 2^2\cdot0.2$
  \item $\sum_{z=0}^2 z^2$
  \end{enumerate}
\item If $E[Z]\approx0.7$, compute $\mathrm{Var}(Z)=E[Z^2]- (E[Z])^2.$
\end{enumerate}

\noindent
\textbf{6.1. Answer}
\begin{itemize}
\item (a) $E[Z]=0\times0.5+1\times0.3+2\times0.2=0.7.$ Correct: \textbf{(B)}.
\item (b) $E[Z^2]=0^2\times0.5 + 1^2\times0.3 + 2^2\times0.2=1.1.$ Correct: \textbf{(C)}.
\item (c) $\mathrm{Var}(Z)=1.1 - (0.7)^2=1.1-0.49=0.61.$
\end{itemize}

\textbf{6.2. Question}\\
In one sentence, why might standard deviation be more straightforward than variance?
\begin{enumerate}[label=\Alph*.]
\item Because SD has units matching the original data
\item Because variance is always negative
\item Because standard deviation is never used
\item Because standard deviation only works with discrete variables
\end{enumerate}

\noindent
\textbf{6.2. Answer:} \textbf{(A)} The standard deviation is in the same units as the original data, making it more intuitive.

\hrule

\subsection*{7. Large Sample Theorems}

\textbf{7.1. Question}\\
The Law of Large Numbers (LLN) states that as $n$ grows large, the sample mean:
\begin{enumerate}[label=\Alph*.]
\item Always moves away from the true mean
\item Converges to the distribution's mode
\item Converges to the distribution's median
\item Converges to the population mean
\end{enumerate}

\noindent
\textbf{7.1. Answer:} \textbf{(D)} Converges to the population mean.

\bigskip

\textbf{7.2. Question}\\
Exam scores have $\mu=70,\,\sigma=10.$ Let $\overline{X}$ be the mean of a random sample of $n=36.$ By the CLT:
\begin{enumerate}[label=(\alph*)]
\item $\overline{X}$ is approximately distributed as:
  \begin{enumerate}[label=\Alph*.]
  \item $\mathrm{Uniform}(70,10)$
  \item $\mathrm{Normal}\!\Bigl(70,\frac{10^2}{36}\Bigr)$
  \item $\mathrm{Normal}(70,10)$
  \item $\mathrm{Binomial}(36,0.7)$
  \end{enumerate}

\item $P(\overline{X}>75)$ is:
  \begin{enumerate}[label=\Alph*.]
  \item $1-\Phi\!\bigl(\tfrac{75-70}{\,10/\sqrt{36}\,}\bigr)$
  \item $1-\Phi\!\bigl(\tfrac{75-70}{10}\bigr)$
  \item $\Phi\!\bigl(\tfrac{70-75}{\,10/\sqrt{36}\,}\bigr)$
  \item $\sum_{k=75}^{100}\binom{36}{k}(0.7)^k(0.3)^{36-k}$
  \end{enumerate}

\item Which sample size makes the CLT approximation more accurate?
  \begin{enumerate}[label=\Alph*.]
  \item $n=36$
  \item $n=100$
  \item Both are the same
  \item No difference
  \end{enumerate}
\end{enumerate}

\noindent
\textbf{7.2. Answer}
\begin{itemize}
\item (a) \textbf{(B)} $\mathrm{Normal}\!\bigl(70,\tfrac{10^2}{36}\bigr).$
\item (b) \textbf{(A)} $1-\Phi\!\Bigl(\tfrac{75-70}{\,10/\sqrt{36}\,}\Bigr).$
\item (c) \textbf{(B)} $n=100$ is more accurate.
\end{itemize}

\hrule

\subsection*{8. Putting It All Together}

\textbf{8.1. Question}\\
You want to model the number of emails you receive daily. Which distribution is \emph{least} appropriate?
\begin{enumerate}[label=\Alph*.]
\item Binomial
\item A Poisson-like distribution
\item Normal (if large count, roughly symmetric)
\item $\mathrm{Uniform}[0,1]$
\end{enumerate}

\noindent
\textbf{8.1. Answer:} \textbf{(D)} $\mathrm{Uniform}[0,1]$ is least appropriate because daily emails are nonnegative integers, not a continuous fraction in $[0,1]$.

\bigskip

\textbf{8.2. Question}\\
A student flips a fair coin $10$ times and gets $10$ heads. They say, ``My next flip \textit{must} be tails!'' Which statement best describes this reasoning?
\begin{enumerate}[label=\Alph*.]
\item Correct: next flip is guaranteed tails
\item Incorrect: each flip is independent, so it's still $0.5$
\item Possibly correct: the coin will ``balance out''
\item Impossible to determine
\end{enumerate}

\noindent
\textbf{8.2. Answer:} \textbf{(B)}. Each flip is independent, so probability of tails remains $0.5.$

\end{document}